\emph{Setting.} The derived metric and the exponent sets are those of Theorem \claimref{synthesis.kasner-approach-gives-horizon-exit-and-index-three}(i): the Kasner form near a wall (Lemma \claimref{interior.kasner-form-interior}), with $\sum p = 1$ over all directions at the vacuum points ($\sum p^{2} = 1$) and at the points that stiff matter reaches. Theorem \claimref{synthesis.kasner-wall-matter-maxwell-bounces-stiff-moves-residue} proves the latter in four dimensions (stiff matter keeps $\sum p = 1$ and gives $\sum p^{2} = 1 - c^{2} < 1$); its five-dimensional analogue, with $\sum_{j}P_{j} = 1$ and $\sum_{j}P_{j}^{2} = 1 - c^{2}$ over four spatial directions, is standard (cited from memory) and is not in the store, which holds no five-dimensional stiff source, so the interior points of the five-dimensional ball are a scope extension and the vacuum points are the only ones the store itself reaches. And the hidden direction is the five-dimensional towers' direction $y$ of Proposition \claimref{synthesis.unit-ball-clauses-hold-at-one-epoch-unless-hidden-direction-co-expands}, the derived metric of the diagonal field $\hb^{-1}(e_{y}) = b\,e_{y}$, read here near a wall with $b = \tau^{p_{h}}$.

\emph{(i)} The action's kinetic weight is $W(\tau) = \sqrt{|g_{4}|}\,b^{\lambda} = \tau^{p_{1} + p_{2} + p_{3}}\tau^{\lambda p_{h}} = \tau^{1 - p_{h} + \lambda p_{h}} = \tau^{\alpha}$, using $\sum p = 1$. For a mode $e^{I_{4}kx^{i}}$ (numbers are read in $\mathrm{span}\{1, I_{4}\}$) the field equation $W^{-1}\partial_{\tau}(W\dot\phi) + k^{2}\tau^{-2p_{i}}\phi = 0$ is the stated one. The equation $y'' + \tfrac{1 - 2a}{x}y' + c^{2}x^{2c - 2}y = 0$ has the solutions $x^{a}Z_{a/c}(x^{c})$, $Z$ any Bessel function of order $|a|/c$ (Fact \claimref{synthesis.imported-bessel-reduction-and-small-argument-hankel-forms}); here, after scaling $\tau$ so that the constant is absorbed into $s = k\tau^{q}/q$, $1 - 2a = \alpha$ and $c = q$, so $a = (1 - \alpha)/2 = -(\lambda - 1)p_{h}/2$ and $\nu = |a|/q = |\lambda - 1||p_{h}|/(2q)$. The check verifies the operator identity
\[
\ddot\phi + \frac{\alpha}{\tau}\dot\phi + k^{2}\tau^{-2p_{i}}\phi = k^{2}\tau^{-2p_{i}}\tau^{a}\Bigl(Z_{ss} + \frac{Z_{s}}{s} + \Bigl(1 - \frac{\nu^{2}}{s^{2}}\Bigr)Z\Bigr), \qquad \phi = \tau^{a}Z(s),
\]
exactly on the terms $Z = s^{n}(\ln s)^{m}$, $n \in \{0, 1, 2, \tfrac52, \nu\}$, $m \le 1$, at $k = q$ (the identity is homogeneous in $k$), for $\lambda \in \{-1, 0, 1, 3, 5\}$ at six exponent sets with rational exponents, and that order $0$ fails unless $(\lambda - 1)p_{h} = 0$.

\emph{(ii)} The momentum conjugate to $\phi$ is $W\dot\phi$, so the normalisation is $\tau^{\alpha}(\phi\dot\phi^{*} - \phi^{*}\dot\phi) = I_{4}$. For $\phi = N\tau^{a}H^{(1)}_{\nu}(s)$ the left side is $|N|^{2}\tau^{\alpha + 2a}(ds/d\tau)W[H^{(1)}_{\nu}, H^{(2)}_{\nu}](s)$, and with $\alpha + 2a = 1$, $ds/d\tau = qs/\tau$ and $W = -4I_{4}/(\pi s)$ it is $-4I_{4} q|N|^{2}/\pi$: $|N|^{2} = \pi/(4q)$, free of $k$. Outside ($s \to 0$), $|H^{(1)}_{\nu}(s)|^{2} \sim \Gamma(\nu)^{2}(2/s)^{2\nu}/\pi^{2}$ for $\nu > 0$ (Fact \claimref{synthesis.imported-bessel-reduction-and-small-argument-hankel-forms}), so at fixed $\tau$, with $s \propto k$, $\Delta^{2}(k) = k^{3}|\phi_{k}|^{2}/(2\pi^{2}) \propto k^{3 - 2\nu}$ and $n - 1 = 3 - 2\nu$. The check verifies $\alpha + 2a = 1$ at every listed point and coupling and pins $(\nu, n - 1)$ at both isotropic vacuum points.

\emph{(iii)} A single exponent of the ball $\sum P = 1$, $\sum P^{2} \le 1$ over four directions lies in $[-\tfrac12, 1]$: at fixed $p_{h}$ the least value of the other three squares is $(1 - p_{h})^{2}/3$ (equal values), so $p_{h}^{2} + (1 - p_{h})^{2}/3 \le 1$, that is $4p_{h}^{2} - 2p_{h} - 2 = 2(2p_{h} + 1)(p_{h} - 1) \le 0$. For $|\lambda - 1| \le 2$, $\nu_{i} = |\lambda - 1||p_{h}|/(2(1 - p_{i})) \le |p_{h}|/(1 - p_{i})$. If $\nu_{i} > 1$ on all three observed axes, then $|p_{h}| > 1 - p_{i}$ for $i = 1, 2, 3$ and, summing with $p_{1} + p_{2} + p_{3} = 1 - p_{h}$, $3|p_{h}| > 2 + p_{h}$, so $p_{h} > 1$ (if $p_{h} \ge 0$) or $p_{h} < -\tfrac12$ (if $p_{h} < 0$), both outside the range. Hence at every point of the ball some observed axis has $\nu_{i} \le 1$; an axis with $p_{i} = 1$ forces the other exponents to vanish, so then $p_{h} = 0$ and every exiting axis has order $0$. Along such an axis the index is at least $1$, so for these couplings no Kasner point gives a spectrum red in every observed direction, without any appeal to the isotropy of the record. The check verifies the factorisation and the range, and samples the ball on a grid of step $\tfrac{1}{12}$ in the observed exponents at $\lambda = -1$ (the extreme case $|\lambda - 1| = 2$): the largest over points of the least order over the three axes is $1$, attained.

\emph{(iv)} With $p_{1} = p_{2} = p_{3} = p$, $\sum p = 1$ gives $p_{h} = 1 - 3p$ and $\sum p^{2} = 3p^{2} + (1 - 3p)^{2} = 1 + 12p^{2} - 6p$, which equals $1$ exactly at $p \in \{0, \tfrac12\}$ (the Milne point and $(\tfrac12, \tfrac12, \tfrac12; -\tfrac12)$) and is at most $1$ exactly on $0 \le p \le \tfrac12$ (the interior being five-dimensional stiff points). There $\nu = \tfrac{|\lambda - 1|}{2}\,\tfrac{|1 - 3p|}{1 - p}$, and $1 - \tfrac{|1 - 3p|}{1 - p}$ equals $\tfrac{2p}{1 - p} \ge 0$ for $p \le \tfrac13$ and $\tfrac{2 - 4p}{1 - p} \ge 0$ for $\tfrac13 \le p \le \tfrac12$, vanishing only at $p = 0$ and $p = \tfrac12$. Hence $\nu \le |\lambda - 1|/2$. If the observed exponents are not all equal and $(\lambda - 1)p_{h} \ne 0$, $\nu$ depends on the axis through $1 - p_{i}$ (the check's example $(\tfrac{5}{12}, \tfrac{5}{12}, \tfrac14; -\tfrac{1}{12})$ at $\lambda = -1$ gives $\tfrac17$ and $\tfrac19$), so the index differs between directions; the record's statistical isotropy (a quadrupolar modulation of the primordial spectrum is bounded at the percent level by Planck, Fact \claimref{synthesis.imported-planck-2018-isocurvature-and-isotropy-bounds}) excludes an order-one difference over the observed range; (iii) does not need this. With $p_{h} = 0$ every coupling gives order $0$.

\emph{(v)} Take the Clifford algebra $\Cl(1,4)$ with $e_{0}^{2} = 1$, $e_{j}^{2} = -1$, and the frame $f_{0} = e_{0}$, $f_{j} = \tau^{p_{j}}e_{j}$, $f_{y} = b\,e_{4}$. The product of the orthogonal frame vectors is $\tau^{\sum p}$ times the pseudoscalar, so $\sqrt{|g_{5}|} = \sqrt{|g_{4}|}\,b$, and $Q(f_{y}) = -b^{2}$, so the $y$ entry of the inverse derived metric is $-b^{-2}$. A massless five-dimensional scalar's zero mode has kinetic weight $\sqrt{|g_{5}|}$, $\lambda = 1$ (a generic entry: the store holds no five-dimensional scalar field); the five-dimensional gauge field of the towers (Proposition \claimref{dark-tower-matter.five-dimensional-towers}, a $\mathfrak{u}(3)$ gauge theory in five dimensions) has $F_{\mu y}F^{\mu y}$ weighted by $\sqrt{|g_{5}|}\,b^{-2}$, so its component $A_{y}$, the Wilson-line scalar, has $\lambda = -1$; a field of our four-dimensional tower, confined to the observed directions, has weight $\sqrt{|g_{4}|}$, $\lambda = 0$. The towers' other content is the Dirac species of Proposition \claimref{dark-tower-matter.five-dimensional-content}, fermions, whose quasi-free deposits are bounded by Theorem \claimref{synthesis.quasi-free-deposit-is-vacuum-plus-shot-noise}; the hidden direction's own field $b$ is not a spectator, and its perturbation is the curvature perturbation of the stiff contraction it drives in the four-dimensional Einstein frame, of order $0$ (argued, the pre-big-bang curvature spectrum, cited from memory). The check computes the weights with the store's Clifford library at $\tau = 2^{12}$ for six exponent sets. The towers' gauge field also has four-dimensional vector zero modes with weight $b$ (the Kaluza-Klein ansatz, Fact \claimref{synthesis.imported-kaluza-klein-ansatz}); for a transverse component $A_{j}$ with $k$ along $i$ the weight of the kinetic term is $\sqrt{|g_{4}|}\,b\,g^{\tau\tau}g^{jj} \propto \tau^{1 - 2p_{j}}$, so $\alpha = 1 - 2p_{j}$ and the order is $|2p_{j}|/(2(1 - p_{i})) = p/(1 - p) \le 1$ at the equal-exponent points ($\tfrac12$ at the static point $p = \tfrac13$, the conformal value, as a control), and the index of the potential is at least $1$; the field strength carries two more powers of $k$ and is bluer. So the stored scalars have $|\lambda - 1| \le 2$, (iii) applies at every point, and $\nu \le 1$ and $n - 1 \ge 1$ on every axis at every point with equal observed exponents; against $n_{s} - 1 = -0.0351 \pm 0.0042$ (Fact \claimref{synthesis.imported-planck-2018-lcdm-parameters}) the gap exceeds $1.0351$, more than $246$ standard deviations. For a coupling with $|\lambda - 1| > 3$ the index $3 - |\lambda - 1|\,|1 - 3p|/(1 - p)$ is continuous on $[0, \tfrac12]$ and takes every value between $3 - |\lambda - 1|$ and $3$ (the check samples $\lambda = 5$, range $[-1, 3]$), so a red value would be fixed by choosing $p$, one number for one target. The metric's off-diagonal modes (the Kaluza-Klein vector's field strength is blue) and curvature couplings (only a fitted curvature coupling reaches the measured tilt) are treated in a result \claimref{synthesis.kasner-approach-coupled-modes-reach-measured-tilt-only-by-a-fitted-curvature-coupling}, and a local conversion of a Gaussian field with these indices cannot lower them (\claimref{synthesis.local-conversion-cannot-redden-a-gaussian-spectator}).
